\majorbreak
\appendix
\section{Reproducibility and artifact guide}

The archive deliberately separates original material, review evidence,
uncompiled candidates, and this report. Start with the PDF or its TeX source,
then use the following supporting material when selecting an implementation
pilot. None of the candidate files was applied to the source snapshot.

\begin{description}
\item[\pathcode{report/}]
This report's TeX sources, PDF, build script, and typesetting requirements.
The report is editable; \code{bash build.sh} rebuilds it with XeLaTeX. Its
successful PDF build is independent of the unavailable Lean replay.
\item[\pathcode{sources/}]
The inspected, unmodified repository material. This is a focused review
snapshot, not a replacement for a full repository checkout.
\item[\pathcode{provenance/}]
The selected-source manifest, source-fetch records, SHA checks, and scripts
for materializing and authenticating the snapshot.
\item[\pathcode{inventory/}]
Per-module source metrics, declaration-span triage, import edges, scope-aware
totals, and the inventory narrative. CSV data are suitable for sorting;
JSON records preserve the measurement definitions.
\item[\pathcode{tactics/}]
Local proof candidates, source replacement excerpts, pinned upstream API
evidence, and source-diff measurements. The portable experiment imports
\code{Mathlib}; that broad import is not a production recommendation.
\item[\pathcode{architecture/}]
Abstraction and dependency proposals, coordinate-packaging candidates,
compatibility tests, and the static import-graph/repetition evidence.
\item[\pathcode{circuits/}]
The two indexed circuit facade candidates, extracted exponent labels,
generated-proof census, static arithmetic checks, and analysis script.
\item[\pathcode{validation/}]
The comment-aware static audit, endpoint list and uncompiled audit module,
axiom-log checker, validation tests, and replay instructions. These helpers
are not substitutes for the actual Lean compiler and kernel.
\end{description}

\subsection{Three measurements that must stay distinct}

\textbf{Physical source size} counts newline-delimited text. A generated list
can be wrapped in many ways without changing its information content, so a
line reduction should be accompanied by byte counts and an explanation of
which old names or data it retains. \textbf{Maintenance surface} concerns how
many independently maintained arguments, public declarations, and duplicated
normal forms exist. It can improve even when compatibility wrappers prevent
a net line reduction. \textbf{Kernel and elaborator cost} must be measured by
building both revisions in the same pinned environment; this review did not
measure it.

The static scripts are explicit about scope. Lexical token counts suppress
comments and strings but do not parse Lean. Import reachability follows only
the selected local source graph; it does not model Lake's glob selection or
external modules. Certificate arithmetic checks validate finite numerical
relations independently but do not establish that the corresponding Lean
declaration elaborates or uses only the intended axioms.

\subsection{Reference conventions}

All source links in the report are immutable GitHub blob links at the full
commit above. Unless a path is explicitly rooted elsewhere, file references
are relative to
\pathcode{Computability/DiophantineConsistency/Lean/DiophantineConsistency/}.
Line ranges refer to the original snapshot, not to candidate files.
Supporting Markdown reviews contain additional anchors and full reasoning;
the main report selects the changes with the clearest maintenance payoff.

\textbf{Bottom line.} Preserve the literal objects and theorem contracts,
reuse the project's existing semantic infrastructure, and remove repeated
presentation proofs in small replayable steps. A shorter source file is a
useful result only when the same theorem still checks under the same trust
boundary and remains easier to maintain.
